\documentclass{beamer}
\usetheme{Boadilla}

\usepackage{pgfpages}

\setbeamertemplate{footline}[frame number]
\setbeamertemplate{navigation symbols}{} 

\setbeameroption{show notes on second screen}
\pdfinfo{/Keywords (SP-Right)}

\setbeamerfont{note page}{size=\Large}
\setbeamertemplate{note page}[plain]

\setbeamertemplate{itemize item}[circle]

\setbeamertemplate{sections/subsections in toc}[circle]
\setbeamertemplate{blocks}[rounded][shadow=false]
\setbeamertemplate{enumerate item}[circle]

\setbeamertemplate{theorems}[numbered]

\usepackage{amsmath,amssymb}
\usepackage{mathtools}
\usepackage{mathpartir}

\newtheorem{proposition}[theorem]{Proposition}
\newtheorem*{defcont}{Definition \arabic{theorem} (cont.)}

\newcommand{\concept}[1]{\textbf{#1}}

% Items

\newcommand{\It}[1]{\ensuremath{\mathsf{It}_{#1}}}
\newcommand{\cfitem}[3]{\ensuremath{{[#1 \to #2 \:\cdot\: #3]}}}

% CFGs

% derivation arrows (no prods)

\newcommand{\derivea}{\ensuremath{\Rightarrow[]{}}}
\newcommand{\derivesa}{\ensuremath{\xRightarrow[]{*}}}
\newcommand{\derivena}[1]{\ensuremath{\xRightarrow[]{#1}}}

\newcommand{\derivera}{\ensuremath{\xRightarrow[\mathrm{rm}]{}}}
\newcommand{\deriversa}{\ensuremath{\xRightarrow[\mathrm{rm}]{*}}}
\newcommand{\deriverna}[1]{\ensuremath{\xRightarrow[\mathrm{rm}]{#1}}}

% derivation arrows with prods

\newcommand{\derive}[2]{\ensuremath{{#1} \derivea {#2}}}
\newcommand{\derives}[2]{\ensuremath{{#1} \derivesa {#2}}}
\newcommand{\deriven}[3]{\ensuremath{{#1} \derivena{#2} {#3}}}

\newcommand{\deriver}[2]{\ensuremath{{#1} \derivera {#2}}}
\newcommand{\derivers}[2]{\ensuremath{{#1} \deriversa {#2}}}
\newcommand{\derivern}[3]{\ensuremath{{#1} \deriverna{#2} {#3}}}

% Alphabets

\newcommand{\sym}{\ensuremath{T \cup N}}
\newcommand{\syms}{\ensuremath{{(T \cup N)}^*}}
\newcommand{\syme}{\ensuremath{T \cup N'}}
\newcommand{\symse}{\ensuremath{{(T \cup N')}^*}}


% Automata

\newcommand{\ipda}{\ensuremath{\mathsf{{I}_G}}}
\newcommand{\charfa}{\ensuremath{\mathsf{char(G)}}}
\newcommand{\lrdfa}{\ensuremath{\mathsf{LR_0(G)}}}
\newcommand{\parser}{\ensuremath{\mathsf{P_0(G)}}}
\newcommand{\srpda}{\ensuremath{\mathsf{M_G}}}

% RM Chains and stack words

\newcommand{\chain}[4]{\ensuremath{{#1} \vdash {#2} -\!\!{#3}\!\rightarrow^{*}_{r} #4}}
\newcommand{\sword}[3]{\ensuremath{{#1} \vDash {#2} \vdash_P^* {#3}}}

% Others

\DeclareMathOperator{\completes}{completes}
\DeclareMathOperator{\hist}{hist}
\DeclareMathOperator{\rev}{rev}


\title{Formal Verification of $LR(0)$ Parsing Theory}
\author[Felipe Escall\'{o}n]{Felipe Escall\'{o}n Tapias}
\institute[TUM]{Technical University of Munich}
\date{October 6th, 2026}

\AtBeginSection[]{
  \begin{frame}
    \tableofcontents[currentsection,hideallsubsections]
  \end{frame}
}

\AtBeginSubsection[]{
  \begin{frame}
    \tableofcontents[currentsection,currentsubsection,
      sectionstyle=show/hide,
      subsectionstyle=show/shaded/hide]
  \end{frame}
}

\begin{document}

\begin{frame}
  \titlepage
\end{frame}

\begin{frame}
\frametitle{Outline}
\tableofcontents[hideallsubsections]
\end{frame}

\section{Introduction}

\begin{frame}
\frametitle{Scope of this Thesis}
\begin{itemize}
  \item Formalization $LR(0)$ parsing theory from Wilhelm et al. in 
    \textit{Compiler Design: Syntactic and Semantic Analysis}
  \item Using the Isabelle Proof Assistant
  \begin{itemize}
    \item Building on the formalization of CFG theory by Nipkow et al.
  \end{itemize}
\end{itemize}
\structure{Major Results}
\note{
  \begin{itemize}
    \item Errors in the book corrected
    \item Book showed parser's determinism, correctness remained open
  \end{itemize}
}
\begin{itemize}
  \item Formal verification of the book's $LR(0)$ section
  \item Formal proof of the correctness of the $LR(0)$ parser 
  \begin{itemize}
    \item Not discussed by Wilhelm et al.
  \end{itemize}
\end{itemize}
\end{frame}


\begin{frame}
\frametitle{Notes on Terminology}
\begin{itemize}
\item ``Formal'': machine-checked
\item ``Formalization'': translation of definitions, theorems, and proofs into 
code that is then machine-checked
\end{itemize}
\end{frame}

\section{Basic Definitions}

\subsection{Extended and Reduced Grammars}

\begin{frame}
  \frametitle{Reduced and Extended Grammars}
  \begin{definition}[Reduced grammar]
    A CFG $G$ is \concept{reduced} if its productions contain only useful
    nonterminals.
  \end{definition}
  \begin{definition}[Extended grammar]
    Let $G$ be a CFG with production set $P$ and start symbol $S$.
    The \concept{extension} of $G$ is the CFG resulting from extending $G$ by a
    new start symbol $S'$, and $P$ by the production $S' \rightarrow S$. This resulting
    grammar is therefore called \concept{extended}.
  \end{definition}
\end{frame}

\subsection{Context-Free Items}

\begin{frame}
  \frametitle{Context-Free Items}
  \begin{definition}
    A context-free item for a CFG with productions $P$, terminals $T$ and nonterminals
    $N$ is a triple 
    $(A, \alpha, \beta) \in N \times \syms \times \syms$ such that
    $A \rightarrow \alpha \beta \in P$.\\
    An item $(A, \alpha, \beta)$ is often written as \cfitem{A}{\alpha}{\beta}.
  \end{definition}
  \begin{itemize}
    \item Context-free items allow tracking the current state of the parsing process
    \item Symbols to the right of the bullet are shifted to the left as they are consumed
    \item Empty strings are often written implicitly
    \begin{itemize}
      \item Example: \cfitem{A}{}{\beta} instead of \cfitem{A}{\varepsilon}{\beta}
    \end{itemize}
  \end{itemize}
\end{frame}

\begin{frame}
  \frametitle{Further Definitions}
  \begin{itemize}
    \item An item of the form \cfitem{A}{}{\alpha} is \concept{initial}
    \item An item of the form \cfitem{A}{\alpha}{} is \concept{complete}
    \item The \concept{history} of an item \cfitem{A}{\alpha}{\beta} is $\alpha$
    \item $\hist$: concatenation of the
    histories of the items in an item sequence
    \small
    \[ \hist (\cfitem{A_1}{\alpha_1}{\beta_1}\cfitem{A_2}{\alpha_2}{\beta_2}\ldots\cfitem{A_n}{\alpha_n}{\beta_n})
        =  \alpha_1 \alpha_2 \ldots \alpha_n \]
    \normalsize
    \item \It{G}: set of all context-free items for the CFG $G$
    \pause
    \bigskip
    \item $\rev$: reverse of a sequence
    \[ \rev (i_1 i_2 \ldots i_n) = i_n \ldots i_2 i_1 \]
  \end{itemize}
\end{frame}

\subsection{Generalized Pushdown Automata}

\begin{frame}
  \frametitle{Generalized Pushdown Automata (GPDAs)}
  \note{Delta split into two to match Isabelle implementation}
  \begin{itemize}
    \item Similar to conventional PDAs
    \begin{itemize}
      \item Popping multiple stack symbols possible
    \end{itemize}
    \item GPDA $M$ is a 5-tuple
    \[ M = (\varSigma, Q, q_0, F, \varDelta) \]
    \vspace{-25 pt}
    \begin{itemize}
      \item $\varSigma$: input alphabet
      \item $Q$: set of states
      \item $q_0$: initial state, $q_0 \in Q$
      \item $F$: set of final states, $F \subseteq Q$
      \item $\varDelta$: transition relation
    \end{itemize}
    \vspace{2.5 pt}
    \item \concept{Configuration} for $M$:
    $(\varrho, w) \in Q^+ \times \varSigma^*$
    \begin{itemize}
      \item $\varrho$: current stack 
      \item $w$: remaining input
    \end{itemize}
    \item \concept{Initial} configuration: $(q_0, w)$ for any $w$
    \item \concept{Final} configuration: $(f, \varepsilon)$ for $f \in F$ 
    \item Language of $M$:
    $L(M) = \{w \mid \exists f \in F.\ (q_0, w) \vdash^* (f, \varepsilon)\}$
  \end{itemize}
\end{frame}

\begin{frame}
  \frametitle{GPDA Steps}
  \begin{itemize}
    \item GPDA steps notated by $\vdash$
    \item \concept{computations} = sequences of GPDA steps
    \begin{itemize}
      \item reflexive transitive closure of steps ($\vdash^*$)
    \end{itemize}
  \end{itemize}
\end{frame}

\section{The Item Pushdown Automaton}

\begin{frame}
  \frametitle{Conventions Moving Forward}
  \begin{itemize}
    \item $G$ is a CFG with start symbol $S$ production set $P$, terminals $T$
    and nonterminals $N$
    \item $G'$ is the extension of $G$ with start symbol $S'$, production set $P'$,
    and nonterminals $N'$
  \end{itemize}
\end{frame}

\begin{frame}
  \frametitle{The Item Pushdown Automaton}
  \begin{definition}
    The \concept{item pushdown automaton} (IPDA) to G is the GPDA 
    \[ \ipda = (T, \It{G'}, \cfitem{S'}{}{S}, \{\cfitem{S'}{S}{}\},
      \varDelta_I \cup \varDelta_{I,\varepsilon}) \]
      where
    \begin{align*}
      \varDelta_I &=
    \begin{multlined}[t]
        \{([\cfitem{X}{\beta}{a \gamma}], a, [\cfitem{X}{\beta a}{\gamma}])\\
        \mid X\ \beta\ a\ \gamma.\ (X, \beta a \gamma) \in P'\}
    \end{multlined}\\
      \varDelta_{I,\varepsilon} &= \ldots
    \end{align*}
  \end{definition}
  \begin{itemize}
    \item IPDA steps notated by $\vdash_I$ and $\vdash^{*}_{I}$
  \end{itemize}
\end{frame}

\begin{frame}
  \frametitle{IPDA Transitions}
  \begin{itemize}
    \item \concept{Shifting} transitions ($\varDelta_I$)
    \[ (\cfitem{A}{\alpha}{a \beta} \varrho, aw) \vdash_{I} (\cfitem{A}{\alpha a}{\beta}\varrho, w) \]
    \item \concept{Expanding} transitions ($\varDelta_{I,\varepsilon}$)
    \begin{gather*}
    (\cfitem{A}{\alpha}{B \beta} \varrho, w) \vdash_{I} (\cfitem{B}{}{\gamma}
    \cfitem{A}{\alpha}{B \beta} \varrho, w)\\
    \text{for } B \rightarrow \gamma \in P'
    \end{gather*}
    \item \concept{Reducing} transitions ($\varDelta_{I,\varepsilon}$)
    \[(\cfitem{B}{\gamma}{} \cfitem{A}{\alpha}{B \beta} \varrho, w) \vdash_{I}
    (\cfitem{A}{\alpha B}{\beta} \varrho, w)\] 
  \end{itemize}
\end{frame}

\begin{frame}
  \frametitle{Correctness}
  \begin{theorem}[Correctness of the IPDA]
    $L(I_{G}) = L(G)$
  \end{theorem}
  \begin{itemize}
    \item Proof follows Wilhelm et al.\ closely
    \item Auxiliary lemma strengthened
  \end{itemize}
\end{frame}

\section{Rightmost Chains}

\begin{frame}
  \frametitle{Rightmost Chains}
  \framesubtitle{Motivation}
  \begin{quote}
    Let us now assume that we have a rightmost derivation 
    $S' \deriversa \gamma' A w \derivera \gamma' \alpha \beta$. This derivation
    can be decomposed into
    \begin{align*}
    S' &\derivera \alpha_{1} X_{1} \beta_{1} \deriversa \alpha_{1} X_{1} v_{1}\\
    &\deriversa \ldots \deriversa (\alpha_{1} \ldots \alpha_{n}) X_n (N \ldots v_1)\\
    &\derivera (\alpha_{1} \ldots \alpha_{n}) \alpha \beta (N \ldots v_1)
    \end{align*}
    for $X_n = A$.
  \end{quote}
  \begin{flushright}
  --- Wilhelm et al., p.~107
  \end{flushright}
\end{frame}

\begin{frame}
  \frametitle{Rightmost Chains}
  \framesubtitle{Definition}
  \begin{itemize}
    \item \chain{P}{\alpha}{\varrho}{\beta}: $\alpha$ reaches $\beta$ with
    rightmost chain $\varrho$ under production set $P$
  \end{itemize}
  \begin{gather*}
    \inferrule[refl]{}{\chain{P}{\alpha}{\varepsilon}{\alpha}}\\[10pt]
    \inferrule[step]{
      \chain{P}{\alpha_0}{\varrho}{\alpha X v} \\
      \deriver{\alpha X v}
        {\alpha \alpha' Y \beta v} \\
      \derivers{\beta}{u}
      }{\chain{P}{\alpha_0}{\cfitem{X}{\alpha'}{Y \beta} \varrho}{\alpha \alpha' Y u v}}
  \end{gather*}
  \pause
  \structure{Results}
  \begin{itemize}
    \item An equivalence between rightmost derivations and rightmost chains
    \item Correspondence of IPDA stacks to rightmost chains
  \end{itemize}
\end{frame}

\section{The Characteristic Finite Automaton and Canonical $LR(0)$ Automaton}

\subsection{The Characteristic Finite Automaton}

\begin{frame}
  \frametitle{The Characteristic Finite Automaton}
  \note{Isabelle splits transition function into nxt and eps}
  \begin{definition}
    The \concept{characteristic finite automaton} to G is the NFA \charfa\ with
    \vspace{-3pt}
    \begin{center}
      \begin{tabular}{r l}
        Alphabet: & \syme\\
        States: & \It{G'}\\
        Initial state: & \cfitem{S'}{}{S}\\
        Final states: & $\{\cfitem{A}{\alpha}{\beta} \in \It{G'}
          \mid \beta = \varepsilon\}$\\
        Transition function: & \ldots\\
      \end{tabular}
    \end{center}
  \end{definition}
  \begin{itemize}
    \item \charfa\ steps denoted by $\vdash_c$ and $\vdash_c^*$
  \end{itemize}
\end{frame}

\begin{frame}
  \frametitle{\charfa\ Transitions}
  \begin{itemize}
    \item Reading transitions
    \[ (\cfitem{X}{\alpha}{Y \beta}, Y \delta) \vdash_c (\cfitem{X}{\alpha Y}{\beta}, \delta)\] 
    \item $\varepsilon$-transitions
    \begin{gather*}
    (\cfitem{X}{\alpha}{Y \beta}, \delta) \vdash_c (\cfitem{Y}{}{\gamma}, \delta)\\
    \text{for } Y \rightarrow \gamma \in P'
    \end{gather*}
  \end{itemize}
\end{frame}

\begin{frame}\label{thm341}
  \frametitle{Theorem 3.4.1: the Relation between \charfa, Rightmost Derivations, 
  and \ipda}
  \begin{quote}
    Let $G$ be a CFG and $\gamma \in \symse $. The following three
    statements are equivalent:
    \begin{enumerate}
      \item There exists a computation
      $(\cfitem{S'}{}{S}, \gamma) \vdash_c^* (\cfitem{A}{\alpha}{\beta}, \varepsilon)$
      of the characteristic finite automaton \charfa.
      \item There exists a rightmost derivation 
      \mbox{\derivers{S'}{\gamma' A w} $\derivera \gamma' \alpha \beta w$} with
      $\gamma = \gamma' \alpha$
      \item There exists a computation 
      $(\cfitem{A}{\alpha}{\beta} \varrho, w)
        \vdash_I^* ([\cfitem{S'}{S}{}] ,\varepsilon)$ of the IPDA \ipda\ such 
        that $\gamma = \hist (\rev (\varrho)) \alpha$ holds.
    \end{enumerate}
  \end{quote}
  \begin{itemize}
    \item Errors corrected
    \item Theorem of equivalence weaker, but individual implications unchanged
  \end{itemize}
\end{frame}

\subsection{The Canonical $LR(0)$ Automaton}

\begin{frame}
  \frametitle{The Canonical $LR(0)$ Automaton}
  \begin{definition}[The Canonical $LR(0)$ Automaton]
    The \concept{canonical $LR(0)$ automaton} \lrdfa\ is the
    DFA resulting from the powerset construction of \charfa\ restricted to
    reachable states.
  \end{definition}
\end{frame}

\section{The $LR(0)$ Parser}

\subsection{Definition}

\begin{frame}
  \frametitle{The $LR(0)$ Parser}
  \note{
    \begin{itemize}
      \item States restricted to Q explicitly akin to the IPDA
      \item Extra: X neq S' avoids conflicts between finish and reduce
      \item f is arbitrary item not in Q
    \end{itemize}
  }
  \begin{definition}
    Let $\varDelta_G$ be the transition function of \lrdfa, $q_{0,G}$ the initial
    state of \lrdfa, $Q$ the set of states of \lrdfa, and $f$ the set
    $\{\cfitem{S'}{}{}\}$. The \concept{$LR(0)$ parser} to G is the 
    GPDA 
    \[ \parser = (T, Q \cup \{f\}, q_{0,G}, \{f\}, \varDelta_P \cup \varDelta_{P,\varepsilon})\]
    where
    \begin{align*}
      \varDelta_P &= \{([q], a, [\varDelta_G(q,a), q])
        \mid q \in Q \land \varDelta_G(q,a) \neq \emptyset\}\\
      \varDelta_{P,\varepsilon} &= \ldots
    \end{align*}
  \end{definition}
  \begin{itemize}
    \item Parser steps denoted by $\vdash_P$ and $\vdash_P^*$
  \end{itemize}
\end{frame}

\begin{frame}
  \frametitle{\parser\ Transitions}
  \begin{itemize}
    \item \concept{Reading} transitions ($\varDelta_P$)
    \begin{gather*}
    (q \varrho, aw) \vdash_P (\varDelta_G(q, a) q \varrho, w)\\
    \text{if } \varDelta_G(q, a) \neq \emptyset
    \end{gather*}
    \item \concept{Reducing} transitions ($\varDelta_{P, \varepsilon}$)
    \begin{gather*}
    (q_n \ldots q_1 q \varrho, w) \vdash_P (\varDelta_G(q, X) q \varrho, w)\\
    \text{ if } \cfitem{X}{\alpha}{} \in q_n \text{ with } |\alpha| = n
    \text{ and } X \neq S'
    \end{gather*}
    \item \concept{Finishing} transitions ($\varDelta_{P, \varepsilon}$)
    \begin{gather*}
      (q q_0 \varrho, w) \vdash_P (f \varrho, w)\\
      \text{ if } \cfitem{S'}{S}{} \in q
    \end{gather*}
  \end{itemize}
\end{frame}

\subsection{$LR(0)$ Adequacy}

\begin{frame}
  \frametitle{$LR(0)$ Adequate and Inadequate States}
  \begin{definition}[Shift-reduce conflict]
    A state $q$ of \parser\ has a \concept{shift-reduce conflict} if 
    it allows for \parser\ to perform both a shift and a reduce transition.
  \end{definition}
  \begin{definition}[Reduce-reduce conflict]
    A state $q$ of \parser\ has a \concept{reduce-reduce conflict} if 
    reducing transitions according to two distinct productions are possible in $q$.
  \end{definition}
  \begin{definition}[$LR(0)$-adequacy]
    A state of \parser\ is \concept{$LR(0)$-inadequate} if it has a shift-reduce
    or a reduce-reduce conflict. Otherwise, it is \concept{$LR(0)$-adequate}.
  \end{definition}
  \pause
  \begin{itemize}
    \item Proof that every $LR(0)$-adequate state falls into one of three cases
    \begin{itemize}
      \item Omitted by Wilhelm et al.
    \end{itemize}
  \end{itemize}
\end{frame}

\subsection{$LR(k)$ Grammars}

\begin{frame}
  \frametitle{String Prefixes and Suffixes}
  \begin{itemize}
    \item $w|_n$: prefix of $w$ of length $n$
    \[(a_0 \ldots a_{n-1} a_n \ldots a_m)|_n = a_0 \ldots a_{n-1}\]
    \item ${}_n|w$: suffix of $w$ starting on index $n$ 
    \[{}_n|(a_0 \ldots a_{n-1} a_n \ldots a_m) = a_n \ldots a_{m}\]
  \end{itemize}
\end{frame}

\begin{frame}
  \frametitle{$LR(k)$ Grammars}
  \begin{definition}[$LR(k)$ Grammar]
    $G'$ is an \concept{$LR(k)$ grammar} if 
    \begin{gather*}
      \derivers{S'}{\alpha X w} \derivera \alpha \beta w \text{ and}\\
      \derivers{S'}{\gamma Y x} \derivera \alpha \beta y \text{ and}\\
      w|_k = y|_k
    \end{gather*}
    implies $\alpha = \gamma$, $X = Y$, and $x = y$.
  \end{definition}
  \begin{alertblock}{Remark}
    This definition is restricted to the extension of a grammar.
  \end{alertblock}
\end{frame}

\begin{frame}
  \frametitle{Theorem 3.4.2: Equivalence of the $LR(0)$ Condition and
    $LR(0)$ Adequacy}
  \begin{theorem}\label{lr0_iff_no_inad}
    $G'$ is an $LR(0)$ grammar if and only if \lrdfa\ has no $LR(0)$-inadequate states.
  \end{theorem}
  \begin{itemize}
    \item ``$\implies$'': correct
    \item ``$\impliedby$'': multiple errors
    \begin{itemize}
      \item Fully rewritten
    \end{itemize}
  \end{itemize}
\end{frame}

\begin{frame}
  \frametitle{Preservation of the $LR(k)$ Condition in Extended Grammars}
  \note{
    \begin{itemize}
      \item Grammars with left recursions standard; LR(0) problematic and needs care
      \item Parser nondeterministic for unextended LR(0) grammars with such left recursions
      \item Cycle = ambiguous grammar
      \begin{itemize}
        \item Not used in practice
        \item LR(k+1) not problematic
      \end{itemize}
    \end{itemize}
  }
  \begin{itemize}
    \item Original $LR(k)$ definition for extended grammars only
    \item For CFG $G$ with extension $G'$: $G\ is\ LR(k) \overset{?}{=} G'\ is\ LR(k)$
    does not hold in general
  \end{itemize}
  \begin{theorem}[Preservation of the $LR(0)$ condition]
    $G'$ is an $LR(0)$ grammar if and only if $G$ is an $LR(0)$ grammar with no
    left recursions of the form \derivern{S}{n+1}{S \alpha}.
  \end{theorem}
  \begin{theorem}[Preservation of the $LR(k)$ condition for $k > 0$]
    $G'$ is an $LR(k)$ grammar for $k > 0$ if and only if $G$ is an $LR(k)$
    grammar with no cycles of the form \derivern{S}{n+1}{S}.
  \end{theorem}
\end{frame}

\subsection{Correctness}

\begin{frame}
  \frametitle{Stack Words: Proving Soundness}
  \begin{itemize}
    \item Symbols passed to $\varDelta_G$ tracked as a string (\concept{stack word})
  \pause
    \item \sword{\alpha}{c_0}{c_1}: $c_0$ reaches $c_1$ with stack word $\alpha$
  \end{itemize}
  \begin{gather*}
    \inferrule[refl]{}{\sword{\varepsilon}{c_0}{c_0}}\\[10 pt]
    \inferrule[step]{\sword{\alpha}{c_0}{(\varrho q \sigma, u)} \\
    (\varrho q \sigma, u) \vdash_P (\varDelta_G(q, X) q \sigma, v)}{
      \sword{X ({}_{|\varrho|}|\alpha)}{c_0}{(\varDelta_G(q, X) q \sigma, v)}
    }
  \end{gather*}
  \pause
  \begin{theorem}[Soundness of the $LR(0)$ parser]
    $L(\parser) \subseteq L(G)$
  \end{theorem}
\end{frame}

\begin{frame}
  \frametitle{The Shift-Reduce Pushdown Automaton}
\begin{definition}
  The \concept{shift-reduce pushdown automaton} (SRPDA) for $G$ is the
  GPDA 
  \[ \srpda = (T, \sym \cup \{q_0, q_f\}, q_0, \{q_f\},
    \varDelta_M \cup \varDelta_{M, \varepsilon}) \]
   for $(\sym) \cap \{q_0, q_f\} = \emptyset$ and $\varDelta_M = \ldots$,
   $\varDelta_{M, \varepsilon} = \ldots$
\end{definition}
\begin{itemize}
  \item We denote SRPDA steps by $\vdash_M$ and $\vdash_M^*$
  \item SRPDA definition following lecture slides by Petter
\end{itemize}
\end{frame}

\begin{frame}
  \frametitle{\srpda\ Transitions}
  \begin{itemize}
    \item \concept{Reading} transitions ($\varDelta_M$)
    \[ (\alpha, aw) \vdash_M (a \alpha, w) \]
    \item \concept{Reducing} transitions ($\varDelta_{M, \varepsilon}$)
      \[ (\rev(\alpha)\beta, w) \vdash_M (A \beta, w) \text{ for } (A, \alpha) \in P \]
    \item \concept{Finishing} transitions ($\varDelta_{M, \varepsilon}$)
      \[ (S q_0 \alpha, w) \vdash_M (q_f \alpha, w) \]
  \end{itemize}
\end{frame}

\begin{frame}
  \frametitle{Proving Completeness with the SRPDA}
  \begin{lemma}[Completeness of the SRPDA]
    $L(G) \subseteq L(\srpda)$
  \end{lemma}
  \begin{theorem}[Completeness of the $LR(0)$ parser]
    $L(G) \subseteq L(\parser)$
  \end{theorem}
\end{frame}

\begin{frame}
  \frametitle{Correctness}
  \begin{theorem}[Correctness of the $LR(0)$ parser]
    $L(G) = L(\parser)$
  \end{theorem}
\end{frame}

\section{Conclusion}

\begin{frame}
  \frametitle{Results}
  \begin{itemize}
    \item Reduced grammars, grammar extensions, and GPDAs
    \item IPDA and correctness proof 
    \item Rightmost chains
    \item \charfa\ and correction of Theorem~\hyperref[thm341]{3.4.1}
    \begin{itemize}
      \item Definition of \lrdfa
    \end{itemize}
    \item The $LR(0)$ parser, $LR(0)$-adequacy and proof of $LR(0)$-adequate
    cases
    \begin{itemize}
      \item Small corrections in the definition
    \end{itemize}
    \item $LR(k)$ and correction of Theorem~\ref{lr0_iff_no_inad}
    \item Preservation of $LR(k)$ property under grammar extension
    \begin{itemize}
      \item Necessary and sufficient conditions
    \end{itemize}
    \item Correctness of \parser
    \begin{itemize}
      \item Stack words
      \item SRPDA
    \end{itemize}
  \end{itemize}
\end{frame}

\begin{frame}
  \frametitle{Future Work}
  \note{At final point: alphabet finiteness problematic since it's a type
  \begin{itemize}
  \item finiteness of \It\ important in my formalization
  \end{itemize}}
  \begin{itemize}
    \item Formalization of the notion of determinism (Wilhelm et al., p. 58)
    \begin{itemize}
      \item Proof that the parser is deterministic iff \lrdfa\ has no 
      $LR(0)$-inadequate states
    \end{itemize}
    \item An executable $LR(0)$ parser
    \item Formalization of an $LR(0)$ parser with output (Wilhelm et al., p. 63) 
    \item Formalization of GPDAs and equivalence with PDAs
    \item Formalization of $LR(k)$ theory for arbitrary $k$
  \end{itemize}
\end{frame}

\end{document}
